\documentclass{article}
\usepackage[margin=1in]{geometry}
\usepackage{amsmath, amssymb}
\usepackage{listings}
\usepackage{xcolor}
\usepackage{hyperref}

\usepackage{tcolorbox}
\tcbuselibrary{listings}

\definecolor{lightbg}{HTML}{F7F7F7}
\definecolor{vscodeblue}{HTML}{0000FF}
\definecolor{vscodered}{HTML}{A31515}
\definecolor{vscodegreen}{HTML}{008000}
\definecolor{vscodegray}{HTML}{808080}

\newtcblisting{jupytercode}{
  colback=lightbg,
  colframe=gray!40,
  listing only,
  left=15pt,
  top=0pt,
  bottom=0pt,
  boxrule=0.5pt,
  arc=3pt,
  listing options={
    language=Python,
    basicstyle=\ttfamily\footnotesize\color{black},
    keywordstyle=\color{vscodeblue}\bfseries,
    commentstyle=\color{vscodegreen}\itshape,
    stringstyle=\color{vscodered},
    numberstyle=\tiny\color{vscodegray},
    numbers=left,
    stepnumber=1,
    showstringspaces=false,
    breaklines=true,
    tabsize=4
  }
}

\title{Numerical Methods: Comprehensive Code and Theory Reference}
\author{Python Implementation Reference}
\date{}

\begin{document}

\maketitle
\tableofcontents
\newpage

%%%%%%%%%%%%%%%%%%%%%%%%%%%%%%%%%%%%%%%%%%%%%%%%
% CHAPTER 1: ROOT FINDING
%%%%%%%%%%%%%%%%%%%%%%%%%%%%%%%%%%%%%%%%%%%%%%%%
\section{Root-Finding Methods}

\subsection{Bisection Method}
\textbf{Theory:} The Bisection method finds a root of $f(x) = 0$ by repeatedly halving a bracketing interval $[a, b]$ where $f(a)f(b) < 0$. It evaluates the midpoint $c = (a+b)/2$ and replaces either $a$ or $b$ with $c$ such that the root remains bracketed. \\
\textbf{Error Bounds \& Convergence:} The error after $n$ iterations is bounded by $|x_n - r| \le \frac{b - a}{2^{n+1}}$. The convergence is \textbf{linear}.

\begin{jupytercode}
import numpy as np

def f(x):
    return (x)**3 + x - 1

a = 0.1
b = 1.0
tol = 10**(-8)

# Bisection Method implementation
if f(a)*f(b) > 0:
    print('choose different points')
else:
    while np.abs(b-a)/2 > tol:
        c = 0.5*(a+b)
        if f(a)*f(c) < 0:
            b = c
        else:
            a = c
            
    print('root = ', c)
    print('f(c) = ', f(c))
\end{jupytercode}
\newpage

\subsection{Newton-Raphson Method}
\textbf{Theory:} Uses the tangent line at the current guess to find the next guess. The iteration formula is $x_{i+1} = x_i - \frac{f(x_i)}{f'(x_i)}$. \\
\textbf{Error Bounds \& Convergence:} The error scales quadratically near a simple root: $\epsilon_{n+1} \approx \frac{f''(r)}{2f'(r)} \epsilon_n^2$. Convergence is \textbf{quadratic}.

\begin{jupytercode}
import numpy as np

def f(x):
    return x**3 + x - 1
    
def fd(x):
    return 3*x**2 + 1

x0 = 0.1
tol = 10**(-8)
val = 1

# Newton-Raphson implementation
if np.abs(f(x0)) < tol:
    print('root =', x0)
else:
    while np.abs(val) > tol:
        x = x0 - f(x0)/fd(x0)
        val = (x - x0)
        x0 = x
        
    print('root = ', x)
\end{jupytercode}
\newpage

\subsection{Secant Method}
\textbf{Theory:} Replaces the derivative in Newton's method with a finite-difference slope passing through the last two iterates. \\
\textbf{Error Bounds \& Convergence:} The error scales as $\epsilon_{n+1} \approx C \epsilon_n^\alpha$, where $\alpha = \frac{1+\sqrt{5}}{2} \approx 1.618$ (the golden ratio). Convergence is \textbf{superlinear}.

\begin{jupytercode}
import numpy as np

def f(x):
    return x**3 + x - 1

x0 = 0.1
x1 = 0.2
tol = 10**(-8)
val = 1

# Secant Method implementation
if np.abs(f(x0)) < tol:
    print('root =', x0)
elif np.abs(f(x1)) < tol:
    print('root =', x1)
else:
    while np.abs(val) > tol:
        fd = (f(x1) - f(x0))/(x1 - x0)
        x = x1 - f(x1)/fd
        val = (x - x1)
        x0 = x1
        x1 = x

    print('root = ', x1)
\end{jupytercode}
\newpage

\subsection{Inverse Quadratic Interpolation (IQI)}
\textbf{Theory:} Fits a sideways quadratic curve (where $x = P(y)$) through the last three iterates. 
$$ P(y) = x_0 \frac{(y - y_1)(y - y_2)}{(y_0 - y_1)(y_0 - y_2)} + x_1 \frac{(y - y_0)(y - y_2)}{(y_1 - y_0)(y_1 - y_2)} + x_2 \frac{(y - y_0)(y - y_1)}{(y_2 - y_0)(y_2 - y_1)} $$
The next root guess is found by evaluating $P(0)$, which simplifies to:
$$ x_3 = x_2 - \frac{r(r - q)(x_2 - x_1) + s(1 - r)(x_2 - x_0)}{(q - 1)(r - 1)(s - 1)} $$ \\
\textbf{Error Bounds \& Convergence:} The convergence is \textbf{superlinear}, with an order of $\approx 1.84$. It requires no derivatives but needs three initial guesses.

\begin{jupytercode}
import numpy as np

def f(x):
    return x**3 + x - 1

def iqi(f, x0, x1, x2, tol=1e-8, max_iter=100):
    for i in range(max_iter):
        y0, y1, y2 = f(x0), f(x1), f(x2)
        
        # Check if we are close enough to zero
        if abs(y2) < tol:
            return x2
            
        # These are the terms for the Lagrange interpolation formula (but sideways!)
        # We calculate the next x guess by plugging y=0 into the sideways parabola
        q = f(x0) / f(x1)
        r = f(x2) / f(x1)
        s = f(x2) / f(x0)
        
        # Formula for the next guess (algebraically simplified)
        x3 = x2 - (r * (r - q) * (x2 - x1) + s * (1 - r) * (x2 - x0)) / ((q - 1) * (r - 1) * (s - 1))
        
        # Shift everything over for the next loop
        x0, x1, x2 = x1, x2, x3
        
    return x2

# Example usage
root = iqi(f, 0.1, 0.15, 0.25)
print("Root is:", root)
\end{jupytercode}
\newpage


%%%%%%%%%%%%%%%%%%%%%%%%%%%%%%%%%%%%%%%%%%%%%%%%
% CHAPTER 2: INTERPOLATION
%%%%%%%%%%%%%%%%%%%%%%%%%%%%%%%%%%%%%%%%%%%%%%%%
\section{Interpolation}

\subsection{Lagrange Interpolation}
\textbf{Theory:} Builds the unique degree-$(n-1)$ interpolating polynomial directly as a weighted sum of Lagrange basis polynomials $L_i(x)$.
\begin{equation*}
P(x) = \sum_{i=0}^{n-1} y_i L_i(x) \quad \text{where} \quad L_i(x) = \prod_{j \neq i} \frac{x - x_j}{x_i - x_j}
\end{equation*}
\textbf{Error Bounds:} The interpolation error is $f(x) - P(x) = \frac{f^{(n)}(\xi)}{n!} \prod_{i=0}^{n-1} (x - x_i)$.

\begin{jupytercode}
import numpy as np
import matplotlib.pyplot as plt

def f(x):
    return 1/(1 + 12*x**2)

# define continuous points for plotting
t = np.arange(-1, 1, 0.01)

# interpolation nodes
x = np.array([-1, -0.99, -0.97, -0.92, -0.8, -0.7, -0.5, -0.25, 0, 0.25, 0.5, 0.7, 0.8, 0.92, 0.97, 0.99, 1])

# init lagrange basis polynomial matrix
l = np.zeros((len(t), len(x)))

# compute basis polynomials automatically
for i in range(len(x)):
    p = np.ones_like(t)
    for j in range(len(x)):
        if i != j:
            p = p * (t - x[j]) / (x[i] - x[j])
    l[:, i] = p

# lagrange interpolating polynomial for f(x)
L = np.zeros_like(t)
for i in range(len(x)):
    L += f(x[i]) * l[:, i]

# plotting
plt.plot(t, f(t), label = 'f(t)')
plt.scatter(x, f(x), color = 'red', marker = '.', label = 'Interpolation nodes')
plt.plot(t, L, label = 'Lagrange Interpolating Polynomial')
plt.grid(True)
plt.title("Lagrange Interpolation [Un-equally spaced]")
plt.legend()
plt.show()
\end{jupytercode}
\newpage

\subsection{Newton's Divided-Difference Interpolation}
\textbf{Theory:} Builds the same unique interpolating polynomial recursively using a divided-difference table.
\textbf{Error Bounds:} Same error as Lagrange interpolation since the resulting polynomial is identical. 

\begin{jupytercode}
import numpy as np
import matplotlib.pyplot as plt

# Fine grid for plotting the true function and interpolation
x = np.linspace(0, np.pi / 2, 100)
f = np.sin(x)

# Interpolation nodes and their function values
t = np.array([0, np.pi/6, 2*np.pi/6, 3*np.pi/6])
y = np.sin(t)

# --- 1. Calculate Newton Divided Difference Coefficients ---
n = len(t)
coef = np.zeros([n, n])
coef[:, 0] = y # first column is just the function values

for j in range(1, n):
    for i in range(n - j):
        coef[i][j] = (coef[i+1][j-1] - coef[i][j-1]) / (t[i+j] - t[i])

# The coefficients for the Newton polynomial are the top row of our table
c = coef[0, :]

# --- 2. Evaluate the Newton Interpolating Polynomial ---
P_x = np.zeros_like(x)
for i in range(len(x)):
    val = c[0]
    prod = 1.0
    for j in range(1, n):
        prod *= (x[i] - t[j-1])
        val += c[j] * prod
    P_x[i] = val

# --- 3. Plotting ---
plt.plot(x, f, label='True function: sin(x)')
plt.scatter(t, y, color='red', marker='x', label='Interpolation nodes')
plt.plot(x, P_x, label='Newton Interpolating Polynomial', linestyle='--')
plt.legend()
plt.title("Newton Divided Difference Interpolation")
plt.grid(True)
plt.show()
\end{jupytercode}
\newpage

\subsection{Natural Cubic Spline Interpolation}
\textbf{Theory:} Fits a piecewise-cubic interpolant ensuring continuous value, slope, and curvature at every interior node. The "Natural" boundary condition specifies that the curvature at the endpoints is zero ($S''(x_1) = S''(x_n) = 0$). \\
\textbf{Error Bounds:} The error bound for cubic splines is $|f(x) - S(x)| \le \frac{5}{384} \max |f^{(4)}(x)| h^4$.

\begin{jupytercode}
import numpy as np 
from matplotlib import pyplot as plt 

def compute_natural_cubic_spline(x, y):
    """
    Computes coefficients for the piecewise cubic polynomials:
    S_i(x) = a_i + b_i*(x - x_i) + c_i*(x - x_i)^2 + d_i*(x - x_i)^3
    """
    x = np.array(x, dtype=float)
    y = np.array(y, dtype=float)
    n = len(x)
    
    # 1. Widths of the intervals
    h = np.diff(x)
    
    # 2. Set up the linear system system: A * c = B
    A = np.zeros((n, n))
    B = np.zeros(n)
    
    # Natural boundary conditions: second derivative at endpoints is 0
    A[0, 0] = 1.0
    A[n-1, n-1] = 1.0
    
    # Fill the tridiagonal matrix based on continuity conditions
    for i in range(1, n-1):
        A[i, i-1] = h[i-1]
        A[i, i] = 2 * (h[i-1] + h[i])
        A[i, i+1] = h[i]
        B[i] = 3 * ((y[i+1] - y[i]) / h[i] - (y[i] - y[i-1]) / h[i-1])
        
    # 3. Solve for 'c' coefficients
    c = np.linalg.solve(A, B)
    
    # 4. Derive 'a', 'b', and 'd' coefficients from 'c'
    a = y[:-1]
    b = np.zeros(n - 1)
    d = np.zeros(n - 1)
    
    for i in range(n - 1):
        b[i] = (y[i+1] - y[i]) / h[i] - h[i] * (2 * c[i] + c[i+1]) / 3
        d[i] = (c[i+1] - c[i]) / (3 * h[i])
        
    return a, b, c[:-1], d

def evaluate_spline(x_query, x_data, a, b, c, d):
    """Evaluates the spline at given query coordinates."""
    y_query = []
    for x_val in x_query:
        # Find which interval the query point falls into
        if x_val <= x_data[0]:
            i = 0
        elif x_val >= x_data[-2]:
            i = len(x_data) - 2
        else:
            # Binary search to locate interval index
            i = np.searchsorted(x_data, x_val) - 1
            
        dx = x_val - x_data[i]
        val = a[i] + b[i]*dx + c[i]*(dx**2) + d[i]*(dx**3)
        y_query.append(val)
        
    return np.array(y_query)

# --- Example Execution & Plotting ---
x_points = np.array([0.0, 1.0, 2.0, 4.0])
y_points = np.array([1.0, 3.0, 2.0, 5.0])

a, b, c, d = compute_natural_cubic_spline(x_points, y_points)
x_dense = np.linspace(-1, 5, 200)
y_dense = evaluate_spline(x_dense, x_points, a, b, c, d)

plt.figure(figsize=(8, 5))
plt.scatter(x_points, y_points, color='red', label='Data Points', zorder=5, marker='.')
plt.plot(x_dense, y_dense, color='blue', label='Custom Cubic Spline')
plt.grid(True)
plt.legend()
plt.show()
\end{jupytercode}
\newpage


%%%%%%%%%%%%%%%%%%%%%%%%%%%%%%%%%%%%%%%%%%%%%%%%
% CHAPTER 3: NUMERICAL INTEGRATION
%%%%%%%%%%%%%%%%%%%%%%%%%%%%%%%%%%%%%%%%%%%%%%%%
\section{Numerical Integration}

\subsection{Trapezoidal Rule}
\textbf{Theory:} Approximates the integral by treating the area under the curve as a series of trapezoids.
$$ \text{Area} = \frac{b - a}{2} (f(a) + f(b)) $$
$$ \int_a^b f(x) dx \approx \frac{h}{2} \left[ f(x_0) + 2 \sum_{i=1}^{m-1} f(x_i) + f(x_m) \right] $$ \\
\textbf{Error Bounds:} 
\begin{itemize}
    \item Local Error (Single Panel): $-\frac{h^3}{12}f''(\xi)$
    \item Global Error (Composite): $-\frac{b-a}{12}h^2 f''(\xi)$
\end{itemize}

\begin{jupytercode}
import numpy as np

def f(x):
    return x**2

def trapezoidal_single(f, a, b):
    h = b - a
    area = h * (f(a) + f(b)) / 2
    return area

def trapezoidal_composite(f, a, b, num_panels):
    h = (b - a) / num_panels
    x = np.linspace(a, b, num_panels + 1)
    
    # Area = h/2 * [f(x_first) + 2*f(x_middles) + f(x_last)]
    # The middle points get counted twice because they touch two trapezoids
    area = (h / 2) * (f(x[0]) + 2 * np.sum(f(x[1:-1])) + f(x[-1]))
    return area

# Example usage
print("Single Trapezoid Area:", trapezoidal_single(f, 0, 1))
print("Composite Trapezoid Area (100 panels):", trapezoidal_composite(f, 0, 1, 100))
\end{jupytercode}
\newpage

\subsection{Simpson's Rule}
\textbf{Theory:} Approximates the integral by fitting a parabola through sets of three points (the endpoints and the midpoint).
$$ \text{Area} = \frac{h}{3} (f(a) + 4f(m) + f(b)) $$
$$ \int_a^b f(x) dx \approx \frac{h}{3} \left[ f(x_0) + 4 \sum_{\text{odd } i} f(x_i) + 2 \sum_{\text{even } i} f(x_i) + f(x_m) \right] $$ \\
\textbf{Error Bounds:} 
\begin{itemize}
    \item Local Error (Single Panel): $-\frac{h^5}{90}f^{(4)}(\xi)$
    \item Global Error (Composite): $-\frac{b-a}{180}h^4 f^{(4)}(\xi)$
\end{itemize}

\begin{jupytercode}
import numpy as np

def f(x):
    return x**2

def simpson_single(f, a, b):
    # A single application requires 3 points: the ends and the exact middle
    midpoint = (a + b) / 2
    h = (b - a) / 2
    
    # Formula uses weights: 1, 4, 1
    area = (h / 3) * (f(a) + 4*f(midpoint) + f(b))
    return area

def simpson_composite(f, a, b, num_panels):
    # Simpson's rule needs an EVEN number of panels
    if num_panels % 2 != 0:
        raise ValueError("Number of panels must be even for Simpson's Rule.")
        
    h = (b - a) / num_panels
    x = np.linspace(a, b, num_panels + 1)
    
    # Weights follow a pattern: 1, 4, 2, 4, 2, ... 4, 1
    y = f(x)
    area = (h / 3) * (y[0] + 4 * np.sum(y[1:-1:2]) + 2 * np.sum(y[2:-1:2]) + y[-1])
    return area

# Example usage
print("Single Simpson Area:", simpson_single(f, 0, 1))
print("Composite Simpson Area (100 panels):", simpson_composite(f, 0, 1, 100))
\end{jupytercode}
\newpage

\subsection{Midpoint Rule}
\textbf{Theory:} Approximates the area with flat rectangles evaluated exactly at the midpoint of each panel. It avoids evaluating the function at the endpoints.
$$ \text{Area} = h \cdot f(m) $$
$$ \int_a^b f(x) dx \approx h \sum_{i=1}^{n} f(m_i) $$ \\
\textbf{Error Bounds:} 
\begin{itemize}
    \item Local Error (Single Panel): $\frac{h^3}{24}f''(\xi)$
    \item Global Error (Composite): $\frac{b-a}{24}h^2 f''(\xi)$
\end{itemize}

\begin{jupytercode}
import numpy as np

def f(x):
    return x**2

def midpoint_single(f, a, b):
    midpoint = (a + b) / 2
    h = b - a
    area = h * f(midpoint)
    return area

def midpoint_composite(f, a, b, num_panels):
    h = (b - a) / num_panels
    
    # Get the exact midpoints of every single panel
    midpoints = np.linspace(a + h/2, b - h/2, num_panels)
    
    # We just sum the height of all the rectangles and multiply by width
    area = h * np.sum(f(midpoints))
    return area

# Example usage
print("Single Midpoint Area:", midpoint_single(f, 0, 1))
print("Composite Midpoint Area (100 panels):", midpoint_composite(f, 0, 1, 100))
\end{jupytercode}
\newpage

\subsection{Gaussian Quadrature}
\textbf{Theory:} Optimally chooses the nodes ($t_i$) and weights ($c_i$) to maximize the exactness of the integral.
$$ \int_{-1}^1 f(t) dt \approx \sum_{i=1}^n c_i f(t_i) $$
$$ \int_a^b f(x) dx \approx \frac{b-a}{2} \sum_{i=1}^n c_i f\left( \frac{b-a}{2}t_i + \frac{b+a}{2} \right) $$ \\
\textbf{Error Bounds:} An $n$-point Gaussian Quadrature rule is completely exact for all polynomials up to degree $2n - 1$. The error for general functions is proportional to $f^{(2n)}(\xi)$.

\begin{jupytercode}
import numpy as np

def f(x):
    return np.exp(-x**2 / 2)

def gaussian_quadrature(f, a, b, n):
    # NumPy has a built-in function to generate the perfect optimal points (t) and their weights (c)
    t, c = np.polynomial.legendre.leggauss(n)
    
    # The points it gives are from -1 to 1. We have to map them to our actual interval [a, b]
    mapped_x = 0.5 * (b - a) * t + 0.5 * (b + a)
    
    # Multiply the weights by our function evaluated at those perfect points
    area = 0.5 * (b - a) * np.sum(c * f(mapped_x))
    return area

# Example usage (Notice how accurate it is with just 4 points!)
print("Gaussian Quadrature (4 points):", gaussian_quadrature(f, -1, 1, 4))
\end{jupytercode}
\newpage


%%%%%%%%%%%%%%%%%%%%%%%%%%%%%%%%%%%%%%%%%%%%%%%%
% CHAPTER 4: LINEAR SYSTEMS & NEWTON FOR SYSTEMS
%%%%%%%%%%%%%%%%%%%%%%%%%%%%%%%%%%%%%%%%%%%%%%%%
\section{Linear Systems and Newton's Method for Systems}

\subsection{Gaussian Elimination}
\textbf{Theory:} Transforms a system $A\mathbf{x} = \mathbf{b}$ into an upper-triangular matrix via forward elimination, and then solves for $\mathbf{x}$ using backward substitution. \\
\textbf{Error Bounds \& Stability:} Vulnerable to round-off error if pivots are close to zero. The residual error heavily depends on the condition number of the matrix $A$.

\begin{jupytercode}
import numpy as np

def gauss_elimination(A, b):
    n = len(b)
    # Forward elimination
    for i in range(n):
        for j in range(i+1, n):
            factor = A[j, i] / A[i, i]
            A[j] -= factor * A[i]
            b[j] -= factor * b[i]
            
    # Back substitution
    x = np.zeros(n)
    for i in range(n-1, -1, -1):
        known_chunk = np.dot(A[i, i+1:], x[i+1:])
        x[i] = (b[i] - known_chunk) / A[i, i]
        
    return x

# Example 3x3 system
A = np.array([[3.0, 2.0, -4.0], [2.0, 3.0, 3.0], [5.0, -3.0, 1.0]])
b = np.array([3.0, 15.0, 14.0])

x = gauss_elimination(A.copy(), b.copy())
print("Solution x:", x)
\end{jupytercode}
\newpage

\subsection{LU Decomposition}
\textbf{Theory:} Factorizes $A = LU$ (Lower and Upper triangular matrices) to solve $A\mathbf{x} = \mathbf{b}$ efficiently for multiple $\mathbf{b}$ vectors. We solve $L\mathbf{y} = \mathbf{b}$ followed by $U\mathbf{x} = \mathbf{y}$.

\begin{jupytercode}
import numpy as np

def lu_decomposition(A):
    n = A.shape[0]
    L = np.eye(n)
    U = A.copy()
    
    for i in range(n):
        for j in range(i+1, n):
            factor = U[j, i] / U[i, i]
            U[j] -= factor * U[i]
            L[j, i] = factor
            
    return L, U

def solve_lu(L, U, b):
    # 1. Forward substitution (Solve L y = b)
    y = np.linalg.solve(L, b)
        
    # 2. Backward substitution (Solve U x = y)
    x = np.linalg.solve(U, y)
        
    return x

A = np.array([[3.0, 2.0, -4.0], [2.0, 3.0, 3.0], [5.0, -3.0, 1.0]])
b = np.array([3.0, 15.0, 14.0])

# Factorize once
L, U = lu_decomposition(A)

# Solve for x
x = solve_lu(L, U, b)

print("L matrix:\n", L)
print("\nU matrix:\n", U)
print("\nVerify L @ U == A:\n", np.allclose(L @ U, A))
print("\nSolution x:", x)
\end{jupytercode}
\newpage

\subsection{PA = LU Decomposition (Partial Pivoting)}
\textbf{Theory:} Implements partial pivoting by swapping rows to ensure the largest possible absolute pivot is used. This prevents division by zero and maintains extreme numerical stability.

\begin{jupytercode}
import numpy as np

def pa_lu_decomposition(A):
    n = A.shape[0]
    L = np.eye(n)
    U = A.copy()
    P = np.eye(n)
    
    for i in range(n):
        # Partial Pivoting: find the row with the largest absolute value in the current column
        max_idx = np.argmax(np.abs(U[i:, i])) + i
        
        # Swap rows in U, P, and the previously calculated parts of L
        if i != max_idx:
            # Swap in U
            U[i], U[max_idx] = U[max_idx].copy(), U[i].copy()
            # Swap in P
            P[i], P[max_idx] = P[max_idx].copy(), P[i].copy()
            # Swap in L (only the columns up to i)
            if i > 0:
                L[i, :i], L[max_idx, :i] = L[max_idx, :i].copy(), L[i, :i].copy()
                
        # Standard LU elimination steps
        for j in range(i+1, n):
            factor = U[j, i] / U[i, i]
            L[j, i] = factor
            U[j] -= factor * U[i]
            
    return P, L, U

def solve_pa_lu(P, L, U, b):
    # 1. Apply permutation to b (We are solving L U x = P b)
    Pb = P @ b
    
    # 2. Forward substitution (Solve L y = Pb)
    y = np.linalg.solve(L, Pb)
        
    # 3. Backward substitution (Solve U x = y)
    x = np.linalg.solve(U, y)
        
    return x

# Example using a matrix that would fail standard LU (0 on diagonal)
A = np.array([[0.0, 2.0, 3.0], [1.0, -1.0, 1.0], [-1.0, 1.0, 0.0]])
b = np.array([4.0, 3.0, 1.0])

P, L, U = pa_lu_decomposition(A)
x = solve_pa_lu(P, L, U, b)

print("Verify P @ A == L @ U:\n", np.allclose(P @ A, L @ U))
print("\nSolution x:", x)
\end{jupytercode}
\newpage

\subsection{Newton's Method for Systems of Nonlinear Equations}
\textbf{Theory:} Solves $\mathbf{F}(\mathbf{x}) = \mathbf{0}$ for multi-variable systems. We repeatedly compute the Jacobian matrix $\mathbf{J}(\mathbf{x})$, find its inverse $\mathbf{J}^{-1}$, and update our guess using the standard inverse formula:
$$ \mathbf{x}_{i+1} = \mathbf{x}_i - [\mathbf{J}(\mathbf{x}_i)]^{-1} \mathbf{F}(\mathbf{x}_i) $$ \\
\textbf{Error Bounds \& Convergence:} Like single-variable Newton's Method, convergence is \textbf{quadratic} near the solution, assuming the Jacobian is non-singular.

\begin{jupytercode}
import numpy as np

# We have two equations we want to solve simultaneously:
# 1) x1 - x0^3 = 0
# 2) x0^2 + x1^2 - 1 = 0
def F(x):
    return np.array([
        x[1] - x[0]**3,
        x[0]**2 + x[1]**2 - 1
    ])

# The Jacobian is a matrix of the partial derivatives of those equations
def Jacobian(x):
    return np.array([
        [-3 * x[0]**2,  1],
        [ 2 * x[0],     2 * x[1]]
    ])

def newton_system(F, Jacobian, initial_guess, tol=1e-6, max_iter=100):
    x = np.array(initial_guess, dtype=float)
    
    for i in range(max_iter):
        J = Jacobian(x)
        fx = F(x)
        
        # Find the inverse of the Jacobian
        J_inv = np.linalg.inv(J)
        
        # Calculate the step: - J_inv * F(x)
        step = -(J_inv @ fx)
        
        # Update x: x = x + step
        x = x + step
        
        # If the step was tiny, we have converged!
        if np.linalg.norm(step) < tol:
            print(f"Converged in {i+1} iterations.")
            return x
            
    return x

# Example usage
initial_guess = [1.0, 2.0]
solution = newton_system(F, Jacobian, initial_guess)
print("Solution:", solution)
\end{jupytercode}

\end{document}
